\documentclass{article}
\usepackage{amsmath,amssymb}
\usepackage{xurl}
\usepackage[colorlinks=true,linkcolor=blue,citecolor=blue,urlcolor=blue]{hyperref}
% Shared commands for the notes in docs/paper-gaps/.
% Every note loads this file via % Shared commands for the notes in docs/paper-gaps/.
% Every note loads this file via % Shared commands for the notes in docs/paper-gaps/.
% Every note loads this file via \input{command}. Project-wide notation
% belongs here, never in per-note redefinitions.

% --- Project configuration ------------------------------------------------
\newcommand{\projectIssueBase}{https://github.com/LionSR/QICLean/issues/}
\newcommand{\projectPrBase}{https://github.com/LionSR/QICLean/pull/}
\newcommand{\projectDocBase}{https://sirui-lu.com/QICLean/docs/}
% ---------------------------------------------------------------------------

\DeclareUnicodeCharacter{2080}{\ifmmode{}_{0}\else\textsubscript{0}\fi}
\DeclareUnicodeCharacter{2081}{\ifmmode{}_{1}\else\textsubscript{1}\fi}
\DeclareUnicodeCharacter{2082}{\ifmmode{}_{2}\else\textsubscript{2}\fi}
\DeclareUnicodeCharacter{2099}{\ifmmode{}_{n}\else\textsubscript{n}\fi}

\newcommand{\Lean}{\textsc{Lean}}

% A formal declaration name, upright sans, breakable at underscores.
\ExplSyntaxOn
\NewDocumentCommand{\leanid}{m}{
  \ifmmode
    \textsf{\detokenize{#1}}
  \else
    \group_begin:
    \urlstyle{sf}
    \tl_set:Nn \l_tmpa_tl {#1}
    \tl_replace_all:Nnn \l_tmpa_tl {\_} {_}
    \exp_args:NV \path \l_tmpa_tl
    \group_end:
  \fi
}
\ExplSyntaxOff

% Traceability links; use in footnotes, not in the running prose.
\newcommand{\ghissue}[1]{\href{\projectIssueBase#1}{GitHub issue \##1}}
\newcommand{\ghpr}[1]{\href{\projectPrBase#1}{PR \##1}}
\newcommand{\doclink}[1]{\href{\projectDocBase#1.html}{\path{#1}}}

% Domain notation, using \providecommand so a note-local refinement stays
% possible.
\providecommand{\tr}{\operatorname{tr}}

% Dirac notation and the identity operator, as in blueprint/src/macros.
\usepackage{dsfont}
\providecommand{\ket}[1]{|#1\rangle}
\providecommand{\bra}[1]{\langle #1|}
\providecommand{\braket}[2]{\langle #1 | #2 \rangle}
\providecommand{\ketbra}[2]{|#1\rangle\!\langle #2|}
\providecommand{\Id}{\mathds{1}}
\providecommand{\one}{\mathds{1}}
\providecommand{\C}{\mathbb{C}}
\providecommand{\MN}[1]{M_{#1}(\C)}
\providecommand{\MD}{M_D(\C)}

% Verdict marker: \gapnote{<kind>}{<status>}. Kind is the policy
% classification (clarification, local-correction, scope-restriction,
% unfaithful, false-source, open-gap); status is open, wip (elimination
% underway), resolved, or historical. Machine-read by the site index;
% prints a verdict line.
\newcommand{\gapnote}[2]{\par\noindent\textbf{Verdict:} #1 (#2).\par}
. Project-wide notation
% belongs here, never in per-note redefinitions.

% --- Project configuration ------------------------------------------------
\newcommand{\projectIssueBase}{https://github.com/LionSR/QICLean/issues/}
\newcommand{\projectPrBase}{https://github.com/LionSR/QICLean/pull/}
\newcommand{\projectDocBase}{https://sirui-lu.com/QICLean/docs/}
% ---------------------------------------------------------------------------

\DeclareUnicodeCharacter{2080}{\ifmmode{}_{0}\else\textsubscript{0}\fi}
\DeclareUnicodeCharacter{2081}{\ifmmode{}_{1}\else\textsubscript{1}\fi}
\DeclareUnicodeCharacter{2082}{\ifmmode{}_{2}\else\textsubscript{2}\fi}
\DeclareUnicodeCharacter{2099}{\ifmmode{}_{n}\else\textsubscript{n}\fi}

\newcommand{\Lean}{\textsc{Lean}}

% A formal declaration name, upright sans, breakable at underscores.
\ExplSyntaxOn
\NewDocumentCommand{\leanid}{m}{
  \ifmmode
    \textsf{\detokenize{#1}}
  \else
    \group_begin:
    \urlstyle{sf}
    \tl_set:Nn \l_tmpa_tl {#1}
    \tl_replace_all:Nnn \l_tmpa_tl {\_} {_}
    \exp_args:NV \path \l_tmpa_tl
    \group_end:
  \fi
}
\ExplSyntaxOff

% Traceability links; use in footnotes, not in the running prose.
\newcommand{\ghissue}[1]{\href{\projectIssueBase#1}{GitHub issue \##1}}
\newcommand{\ghpr}[1]{\href{\projectPrBase#1}{PR \##1}}
\newcommand{\doclink}[1]{\href{\projectDocBase#1.html}{\path{#1}}}

% Domain notation, using \providecommand so a note-local refinement stays
% possible.
\providecommand{\tr}{\operatorname{tr}}

% Dirac notation and the identity operator, as in blueprint/src/macros.
\usepackage{dsfont}
\providecommand{\ket}[1]{|#1\rangle}
\providecommand{\bra}[1]{\langle #1|}
\providecommand{\braket}[2]{\langle #1 | #2 \rangle}
\providecommand{\ketbra}[2]{|#1\rangle\!\langle #2|}
\providecommand{\Id}{\mathds{1}}
\providecommand{\one}{\mathds{1}}
\providecommand{\C}{\mathbb{C}}
\providecommand{\MN}[1]{M_{#1}(\C)}
\providecommand{\MD}{M_D(\C)}

% Verdict marker: \gapnote{<kind>}{<status>}. Kind is the policy
% classification (clarification, local-correction, scope-restriction,
% unfaithful, false-source, open-gap); status is open, wip (elimination
% underway), resolved, or historical. Machine-read by the site index;
% prints a verdict line.
\newcommand{\gapnote}[2]{\par\noindent\textbf{Verdict:} #1 (#2).\par}
. Project-wide notation
% belongs here, never in per-note redefinitions.

% --- Project configuration ------------------------------------------------
\newcommand{\projectIssueBase}{https://github.com/LionSR/QICLean/issues/}
\newcommand{\projectPrBase}{https://github.com/LionSR/QICLean/pull/}
\newcommand{\projectDocBase}{https://sirui-lu.com/QICLean/docs/}
% ---------------------------------------------------------------------------

\DeclareUnicodeCharacter{2080}{\ifmmode{}_{0}\else\textsubscript{0}\fi}
\DeclareUnicodeCharacter{2081}{\ifmmode{}_{1}\else\textsubscript{1}\fi}
\DeclareUnicodeCharacter{2082}{\ifmmode{}_{2}\else\textsubscript{2}\fi}
\DeclareUnicodeCharacter{2099}{\ifmmode{}_{n}\else\textsubscript{n}\fi}

\newcommand{\Lean}{\textsc{Lean}}

% A formal declaration name, upright sans, breakable at underscores.
\ExplSyntaxOn
\NewDocumentCommand{\leanid}{m}{
  \ifmmode
    \textsf{\detokenize{#1}}
  \else
    \group_begin:
    \urlstyle{sf}
    \tl_set:Nn \l_tmpa_tl {#1}
    \tl_replace_all:Nnn \l_tmpa_tl {\_} {_}
    \exp_args:NV \path \l_tmpa_tl
    \group_end:
  \fi
}
\ExplSyntaxOff

% Traceability links; use in footnotes, not in the running prose.
\newcommand{\ghissue}[1]{\href{\projectIssueBase#1}{GitHub issue \##1}}
\newcommand{\ghpr}[1]{\href{\projectPrBase#1}{PR \##1}}
\newcommand{\doclink}[1]{\href{\projectDocBase#1.html}{\path{#1}}}

% Domain notation, using \providecommand so a note-local refinement stays
% possible.
\providecommand{\tr}{\operatorname{tr}}

% Dirac notation and the identity operator, as in blueprint/src/macros.
\usepackage{dsfont}
\providecommand{\ket}[1]{|#1\rangle}
\providecommand{\bra}[1]{\langle #1|}
\providecommand{\braket}[2]{\langle #1 | #2 \rangle}
\providecommand{\ketbra}[2]{|#1\rangle\!\langle #2|}
\providecommand{\Id}{\mathds{1}}
\providecommand{\one}{\mathds{1}}
\providecommand{\C}{\mathbb{C}}
\providecommand{\MN}[1]{M_{#1}(\C)}
\providecommand{\MD}{M_D(\C)}

% Verdict marker: \gapnote{<kind>}{<status>}. Kind is the policy
% classification (clarification, local-correction, scope-restriction,
% unfaithful, false-source, open-gap); status is open, wip (elimination
% underway), resolved, or historical. Machine-read by the site index;
% prints a verdict line.
\newcommand{\gapnote}[2]{\par\noindent\textbf{Verdict:} #1 (#2).\par}

\title{Interpolation and Word Counting in the AKLV Rank Argument}
\author{}
\date{October 21, 2026}
\begin{document}
\maketitle
\gapnote{local-correction}{resolved}

The resolved correction concerns the interpolation wording in the proof of
Lemma~4.2 of Arad, Kitaev, Landau and Vazirani,
\href{https://arxiv.org/abs/1301.1162v1}{arXiv:1301.1162v1},
label \texttt{l:agsp},
\path{References/1301.1162/1dfrustrated34.tex}, lines~626--644.
The corrected exponent-tuple count and a coefficient-rank bound for commuting
gaps are formalized. The full Hamiltonian operator Schmidt-rank theorem remains
a separate assertion.

\section{Interpolation requires a spanning family}

Let $S$ be a finite set of distinct monomials in complex variables, with
$t=|S|$, and write $m_\alpha(z)$ for the monomial indexed by $\alpha\in S$.
There are $t$ evaluation points $z_0,\ldots,z_{t-1}$ and complex numbers
$c_{\alpha r}$, chosen from $S$ alone, such that
\begin{align}
    \sum_{r=0}^{t-1} c_{\alpha r}m_\beta(z_r)
        &= \begin{cases}1,&\alpha=\beta,\\0,&\alpha\ne\beta.\end{cases}
    \label{eq:aklv13_interp_scalar}
\end{align}
Indeed, the vectors $(m_\beta(z))_{\beta\in S}$ span $\mathbb C^S$:
a linear functional annihilating every such vector gives a polynomial
vanishing at every complex assignment, and hence the zero polynomial.
Choosing a basis from these evaluation vectors gives
(\ref{eq:aklv13_interp_scalar}).
This basis lies in $\mathbb C^S$, not in a space of coefficient operators.

For every complex vector space $V$ and every family $(f_\beta)_{\beta\in S}$
in $V$, put $P_f(z)=\sum_{\beta\in S}m_\beta(z)f_\beta$.
For each selected subset $A\subseteq S$, the same points and scalar coefficients give
\begin{align}
    f_\alpha
        &=\sum_{r=0}^{t-1}c_{\alpha r}P_f(z_r),
    \label{eq:aklv13_interp_vector}\\
    \sum_{\alpha\in A}f_\alpha
        &=\sum_{r=0}^{t-1}\left(\sum_{\alpha\in A}c_{\alpha r}\right)P_f(z_r).
    \label{eq:aklv13_interp_sum}
\end{align}
\begin{samepage}
These identities allow zero and dependent coefficient operators. The scalar
and vector recovery identities are proved by
\leanid{MvPolynomial.exists_monomial_evaluation_recovery} and
\leanid{MvPolynomial.monomial_evaluation_recovery_smul}.
For $S=\varnothing$, there are no coefficients to recover and all sums are empty.
\par
\end{samepage}

In the source, $t=\binom{\ell-k+s}{s}$ and the operators
$P_{i,\ell k}(z_r)$ are said to form a basis of the coefficient-operator
space (lines~637--643). They need only span that space.
For example, if $I\ne0$ and $P(z)=(1+z)I$, both nonzero coefficients equal
$I$, so the coefficient space has dimension one although $t=2$.
The homogeneous example $P(z_0,z_1)=(z_0+z_1)I$ has the same defect and
matches the source's fixed-total-degree setting.
Equation~(\ref{eq:aklv13_interp_sum})
proves precisely the conclusion needed at source line~644, without
requiring the coefficient operators to be linearly independent.

\section{The number of collapsed words}

For integers $0\le k\le\ell$, a word with exactly $k$ occurrences of $H_i$
has, after commuting $A$ and $B$ within each intervening gap, the form
\begin{align}
    A^{a_0}B^{b_0}H_i A^{a_1}B^{b_1}H_i\cdots H_i A^{a_k}B^{b_k},
    \qquad \sum_{r=0}^{k}(a_r+b_r)=\ell-k.
    \label{eq:aklv13_word_form}
\end{align}
There are $2k+2$ nonnegative exponents. Their number of possible tuples is
\begin{align}
    \binom{(\ell-k)+(2k+2)-1}{(2k+2)-1}
        &=\binom{\ell+k+1}{2k+1}.
    \label{eq:aklv13_word_count}
\end{align}
The corrected count is proved by
\leanid{Finset.card_piAntidiag_univ_fin_prod_bool}. Combined with the commuting-gap
span and the algebraic rank inequalities, it gives the coefficient-rank bound
\leanid{Matrix.operatorSchmidtRank_coeff_pow_C_add_C_mul_X_le} when the two
gap operators have ranks at most one.

Different tuples may yield the same operator; the count is an upper bound
on the number of collected operator expressions, not an independence claim.

Source lines~651 and~657 instead print $\binom{\ell+k}{2k+1}$.
Already $k=0$, which belongs to the source's range
$0\le k\le\lfloor\ell/s\rfloor$ for every positive $s$, gives $\ell+1$
patterns $A^aB^{\ell-a}$ rather than the printed $\ell$.
At $k=\ell$, the single pattern $H_i^\ell$ gives one rather than zero;
this also covers the constant-power case $\ell=k=0$ when powers are
assembled into a polynomial.

The corrected upper index also appears in Arad, Landau, Vazirani and
Vidick, \href{https://arxiv.org/abs/1602.08828v2}{arXiv:1602.08828v2},
Proposition~\texttt{prop:Cheby-AGSP-eff},
\path{References/1602.08828/local_alg_arxiv_v2.tex}, lines~1324 and~1328:
the two-cut count contains
$\binom{k+j_1+1}{2j_1+1}\binom{k+j_2+1}{2j_2+1}$.
The interpolation identities do not establish the locality estimates,
word-rank bounds, or subsequent rank assembly. Those remain separate
parts of the operator Schmidt rank argument.

\end{document}
